% Source: https://artofproblemsolving.com/community/c3h3191353_2021_bmmt_individual_tiebreaker_round___berkley_mini_math_tournament
% Stable identity: aops-topic-3191353-q03
Place the equal-circle centers at $(-3,0),(3,0)$ and the radius-$2$ center at $(0,4)$. The sought centers lie on the symmetry axis. For the gap circle of radius $r$ and center $(0,y)$, $y=2-r$ and $9+y^2=(r+3)^2$. Therefore $4-4r=6r$, giving $r=\dfrac{2}{5}$. For the enclosing circle of radius $R$ centered at $(0,y)$, the equations are $9+y^2=(R-3)^2$ and $(y-4)^2=(R-2)^2$. Subtracting gives $-8y+7=2R-5$, or $y=\dfrac{6-R}{4}$. Substitution yields $15R^2-84R-36=0$, hence $R=6$ or $-\dfrac{2}{5}$; only $R=6$ is positive. The sum is $\dfrac{32}{5}$, so $m+n=\boxed{37}$.
\begin{asy}
usepackage("amsmath");
size(7cm);
draw(circle((-3,0),3));draw(circle((3,0),3));draw(circle((0,4),2));draw(circle((0,1.6),.4));draw(circle((0,0),6),dashed);label("$3$",(-3,0));label("$3$",(3,0));label("$2$",(0,4));label("$r$",(0,1.6),E);label("$R=6$",(0,-6),S);
\end{asy}
